
\documentclass[a4paper,11pt]{article}
\usepackage[a4paper, margin=8em]{geometry}

% usa i pacchetti per la scrittura in italiano
\usepackage[french,italian]{babel}
\usepackage[T1]{fontenc}
\usepackage[utf8]{inputenc}
\frenchspacing 

% usa i pacchetti per la formattazione matematica
\usepackage{amsmath, amssymb, amsthm, amsfonts}

% usa altri pacchetti
\usepackage{gensymb}
\usepackage{hyperref}
\usepackage{standalone}

\usepackage{colortbl}

\usepackage{xstring}
\usepackage{karnaugh-map}

% imposta il titolo
\title{Electronics for Embedded Systems notes}
\author{Luca Seggiani}
\date{2026}

% spaziatura liste
\usepackage{enumitem}
\setlist[itemize]{noitemsep, topsep=4pt, partopsep=0pt, parsep=2pt, leftmargin=*}
\setlist[enumerate]{noitemsep, topsep=4pt, partopsep=0pt, parsep=2pt, leftmargin=*}

% imposta lo stile
% usa helvetica
\usepackage[scaled]{helvet}
% usa palatino
\usepackage{palatino}
% usa un font monospazio guardabile
\usepackage{lmodern}

\renewcommand{\rmdefault}{ppl}
\renewcommand{\sfdefault}{phv}
\renewcommand{\ttdefault}{lmtt}

% circuiti
\usepackage{circuitikz}
\usetikzlibrary{babel}

% testo cerchiato
\newcommand*\circled[1]{\tikz[baseline=(char.base)]{
            \node[shape=circle,draw,inner sep=2pt] (char) {#1};}}

% disponi il titolo
\makeatletter
\renewcommand{\maketitle} {
	\begin{center} 
		\begin{minipage}[t]{.8\textwidth}
			\textsf{\huge\bfseries \@title} 
		\end{minipage}%
		\begin{minipage}[t]{.2\textwidth}
			\raggedleft \vspace{-1.65em}
			\textsf{\small \@author} \vfill
			\textsf{\small \@date}
		\end{minipage}
		\par
	\end{center}

	\thispagestyle{empty}
	\pagestyle{fancy}
}
\makeatother

% disponi teoremi
\usepackage{tcolorbox}
\newtcolorbox[auto counter, number within=section]{theorem}[2][]{%
	colback=blue!10, 
	colframe=blue!40!black, 
	sharp corners=northwest,
	fonttitle=\sffamily\bfseries, 
	title=Theorem~\thetcbcounter: #2, 
	#1
}

% disponi definizioni
\newtcolorbox[auto counter, number within=section]{definition}[2][]{%
	colback=red!10,
	colframe=red!40!black,
	sharp corners=northwest,
	fonttitle=\sffamily\bfseries,
	title=Definition~\thetcbcounter: #2,
	#1
}

% disponi codice
\usepackage{listings}
\usepackage[table]{xcolor}

\definecolor{codegreen}{rgb}{0,0.6,0}
\definecolor{codegray}{rgb}{0.5,0.5,0.5}
\definecolor{codepurple}{rgb}{0.58,0,0.82}
\definecolor{backcolour}{rgb}{0.95,0.95,0.92}

\lstdefinestyle{codestyle}{
		backgroundcolor=\color{black!5}, 
		commentstyle=\color{codegreen},
		keywordstyle=\bfseries\color{magenta},
		numberstyle=\sffamily\tiny\color{black!60},
		stringstyle=\color{green!50!black},
		basicstyle=\ttfamily\footnotesize,
		breakatwhitespace=false,         
		breaklines=true,                 
		captionpos=b,                    
		keepspaces=true,                 
		numbers=left,                    
		numbersep=5pt,                  
		showspaces=false,                
		showstringspaces=false,
		showtabs=false,                  
		tabsize=2
}

\lstdefinestyle{shellstyle}{
		backgroundcolor=\color{black!5}, 
		basicstyle=\ttfamily\footnotesize\color{black}, 
		commentstyle=\color{black}, 
		keywordstyle=\color{black},
		numberstyle=\color{black!5},
		stringstyle=\color{black}, 
		showspaces=false,
		showstringspaces=false, 
		showtabs=false, 
		tabsize=2, 
		numbers=none, 
		breaklines=true
}

\lstdefinelanguage{x86}{ 
  keywords={AAA, AAD, AAM, AAS, ADC, ADCB, ADCW, ADCL, ADD, ADDB, ADDW, ADDL, AND, ANDB, ANDW, ANDL,
        ARPL, BOUND, BSF, BSFL, BSFW, BSR, BSRL, BSRW, BSWAP, BT, BTC, BTCB, BTCW, BTCL, BTR, 
        BTRB, BTRW, BTRL, BTS, BTSB, BTSW, BTSL, CALL, CBW, CDQ, CLC, CLD, CLI, CLTS, CMC, CMP,
        CMPB, CMPW, CMPL, CMPS, CMPSB, CMPSD, CMPSW, CMPXCHG, CMPXCHGB, CMPXCHGW, CMPXCHGL,
        CMPXCHG8B, CPUID, CWDE, DAA, DAS, DEC, DECB, DECW, DECL, DIV, DIVB, DIVW, DIVL, ENTER,
        HLT, IDIV, IDIVB, IDIVW, IDIVL, IMUL, IMULB, IMULW, IMULL, IN, INB, INW, INL, INC, INCB,
        INCW, INCL, INS, INSB, INSD, INSW, INT, INT3, INTO, INVD, INVLPG, IRET, IRETD, JA, JAE,
        JB, JBE, JC, JCXZ, JE, JECXZ, JG, JGE, JL, JLE, JMP, JNA, JNAE, JNB, JNBE, JNC, JNE, JNG,
        JNGE, JNL, JNLE, JNO, JNP, JNS, JNZ, JO, JP, JPE, JPO, JS, JZ, LAHF, LAR, LCALL, LDS,
        LEA, LEAVE, LES, LFS, LGDT, LGS, LIDT, LMSW, LOCK, LODSB, LODSD, LODSW, LOOP, LOOPE,
        LOOPNE, LSL, LSS, LTR, MOV, MOVB, MOVW, MOVL, MOVSB, MOVSD, MOVSW, MOVSX, MOVSXB,
        MOVSXW, MOVSXL, MOVZX, MOVZXB, MOVZXW, MOVZXL, MUL, MULB, MULW, MULL, NEG, NEGB, NEGW,
        NEGL, NOP, NOT, NOTB, NOTW, NOTL, OR, ORB, ORW, ORL, OUT, OUTB, OUTW, OUTL, OUTSB, OUTSD,
        OUTSW, POP, POPL, POPW, POPB, POPA, POPAD, POPF, POPFD, PUSH, PUSHL, PUSHW, PUSHB, PUSHA, 
        RCLL, RCR, RCRB, RCRW, RCRL, RDMSR, RDPMC, RDTSC, REP, REPE, REPNE, RET, ROL, ROLB, ROLW,
        ROLL, ROR, RORB, RORW, RORL, SAHF, SAL, SALB, SALW, SALL, SAR, SARB, SARW, SARL, SBB,
        SBBB, SBBW, SBBL, SCASB, SCASD, SCASW, SETA, SETAE, SETB, SETBE, SETC, SETE, SETG, SETGE,
        SETL, SETLE, SETNA, SETNAE, SETNB, SETNBE, SETNC, SETNE, SETNG, SETNGE, SETNL, SETNLE,
        SETNO, SETNP, SETNS, SETNZ, SETO, SETP, SETPE, SETPO, SETS, SETZ, SGDT, SHL, SHLB, SHLW,
        SHLL, SHLD, SHR, SHRB, SHRW, SHRL, SHRD, SIDT, SLDT, SMSW, STC, STD, STI, STOSB, STOSD,
        STOSW, STR, SUB, SUBB, SUBW, SUBL, TEST, TESTB, TESTW, TESTL, VERR, VERW, WAIT, WBINVD,
        XADD, XADDB, XADDW, XADDL, XCHG, XCHGB, XCHGW, XCHGL, XLAT, XLATB, XOR, XORB, XORW, XORL},
  keywordstyle=\color{blue}\bfseries,
  ndkeywordstyle=\color{darkgray}\bfseries,
  identifierstyle=\color{black},
  sensitive=false,
  comment=[l]{\#},
  morecomment=[s]{/*}{*/},
  commentstyle=\color{purple}\ttfamily,
  stringstyle=\color{red}\ttfamily,
  morestring=[b]',
  morestring=[b]"
}

\lstdefinelanguage{riscv}{
  keywords={
    LUI, AUIPC,
    JAL, JALR,
    BEQ, BNE, BLT, BGE, BLTU, BGEU,
    LB, LH, LW, LBU, LHU, LD,
    SB, SH, SW, SD,
    ADDI, SLTI, SLTIU, XORI, ORI, ANDI,
    SLLI, SRLI, SRAI,
    ADD, SUB, SLL, SLT, SLTU, XOR, SRL, SRA, OR, AND,
    ADDIW, SLLIW, SRLIW, SRAIW,
    ADDW, SUBW, SLLW, SRLW, SRAW,
    MUL, MULH, MULHSU, MULHU, DIV, DIVU, REM, REMU,
    MULW, DIVW, DIVUW, REMW, REMUW,
    LR.W, SC.W, AMOSWAP.W, AMOADD.W, AMOXOR.W,
    AMOAND.W, AMOOR.W, AMOMIN.W, AMOMAX.W,
    AMOMINU.W, AMOMAXU.W,
    LR.D, SC.D, AMOSWAP.D, AMOADD.D, AMOXOR.D,
    AMOAND.D, AMOOR.D, AMOMIN.D, AMOMAX.D,
    AMOMINU.D, AMOMAXU.D,
    FLW, FSW,
    FMADD.S, FMSUB.S, FNMSUB.S, FNMADD.S,
    FADD.S, FSUB.S, FMUL.S, FDIV.S, FSQRT.S,
    FSGNJ.S, FSGNJN.S, FSGNJX.S,
    FMIN.S, FMAX.S,
    FCVT.W.S, FCVT.WU.S, FCVT.S.W, FCVT.S.WU,
    FMV.X.W, FMV.W.X,
    FEQ.S, FLT.S, FLE.S,
    FLD, FSD,
    FMADD.D, FMSUB.D, FNMSUB.D, FNMADD.D,
    FADD.D, FSUB.D, FMUL.D, FDIV.D, FSQRT.D,
    FSGNJ.D, FSGNJN.D, FSGNJX.D,
    FMIN.D, FMAX.D,
    FCVT.W.D, FCVT.WU.D, FCVT.D.W, FCVT.D.WU,
    FCVT.S.D, FCVT.D.S,
    FMV.X.D, FMV.D.X,
    FEQ.D, FLT.D, FLE.D,
    CSRRW, CSRRS, CSRRC,
    CSRRWI, CSRRSI, CSRRCI,
    FENCE.I,
    FENCE,
    ECALL, EBREAK,
    MRET, SRET, URET,
    WFI,
    SFENCE.VMA,
    NOP, LI, LA, MV,
    NOT, NEG, NEGW,
    SEQZ, SNEZ, SLTZ, SGTZ,
    BEQZ, BNEZ, BLEZ, BGEZ, BLTZ, BGTZ,
    BGT, BLE, BGTU, BLEU,
    J, JR, RET,
    CALL, TAIL,
    RDINSTRET, RDCYCLE, RDTIME,
    CSRR, CSRW, CSRS, CSRC,
    CSRWI, CSRSI, CSRCI
  },
  keywordstyle=\color{blue}\bfseries,
  ndkeywordstyle=\color{darkgray}\bfseries,
  identifierstyle=\color{black},
  sensitive=false,
  comment=[l]{\#},
  morecomment=[s]{/*}{*/},
  commentstyle=\color{purple}\ttfamily,
  stringstyle=\color{red}\ttfamily,
  morestring=[b]',
  morestring=[b]"
}

\lstdefinelanguage{arm}{
  keywords={
    ADC, ADD, AND, BIC, CMN, CMP, EOR, MOV, MVN,
    RSB, RSC, SBC, SUB, TST, TEQ,
    MLA, MLS, MUL, SMLAL, SMULL, UMLAL, UMULL,
    QADD, QSUB, QDADD, QDSUB,
    QADD16, QADD8, QASX, QSAX,
    QSUB16, QSUB8, QSAX, QSUBADDX,
    ASR, LSL, LSR, ROR, RRX,
    B, BL, BX, BLX,
    BXJ,
    BEQ, BNE, BCS, BHS, BCC, BLO,
    BMI, BPL, BVS, BVC,
    BHI, BLS, BGE, BLT, BGT, BLE,
    BAL,
    LDR, LDRB, LDRH, LDRSB, LDRSH,
    LDRD, LDM, LDMIA, LDMIB, LDMDA, LDMDB,
    LDRT, LDRBT, LDRHT, LDRSBT, LDRSHT,
    STR, STRB, STRH, STRD,
    STM, STMIA, STMIB, STMDA, STMDB,
    STRT, STRBT, STRHT,
    PUSH, POP,
    SWP, SWPB,
    REV, REV16, REVSH,
    BFC, BFI, BFXIL, SBFX, UBFX,
    CLZ,
    LDREX, LDREXB, LDREXH, LDREXD,
    STREX, STREXB, STREXH, STREXD,
    DMB, DSB, ISB,
    MCR, MCRR, MRC, MRRC,
    MRS, MSR,
    CPS, SETEND, SVC, BKPT,
    CLREX, WFE, WFI, SEV,
    YIELD, NOP,
    VLDR, VSTR,
    VMOV, VABS, VNEG,
    VADD, VSUB, VMUL, VDIV,
    VMLA, VMLS, VNMLA, VNMLS, VNMUL,
    VMAX, VMIN,
    VSQRT,
    VCMP, VCMPE,
    VCVT,
    VCLZ, VCLS,
    VAND, VORR, VEOR, VBIC, VMVN,
    VLD1, VLD2, VLD3, VLD4,
    VST1, VST2, VST3, VST4,
    ADR, ADCS, ADDS, SUBS, MOVS, ANDS, ORRS, EORS
  },
  keywordstyle=\color{blue}\bfseries,
  ndkeywordstyle=\color{darkgray}\bfseries,
  identifierstyle=\color{black},
  sensitive=false,
  comment=[l]{@},
  morecomment=[l]{;},
  morecomment=[s]{/*}{*/},
  commentstyle=\color{purple}\ttfamily,
  stringstyle=\color{red}\ttfamily,
  morestring=[b]',
  morestring=[b]"
}

\lstset{style=codestyle, language=C}

% disponi sezioni
\usepackage{titlesec}

\titleformat{\section}
	{\sffamily\Large\bfseries} 
	{\thesection}{1em}{} 
\titleformat{\subsection}
	{\sffamily\large\bfseries}   
	{\thesubsection}{1em}{} 
\titleformat{\subsubsection}
	{\sffamily\normalsize\bfseries} 
	{\thesubsubsection}{1em}{}

% tikz
\usepackage{tikz}

% float
\usepackage{float}

% grafici
\usepackage{pgfplots}
\pgfplotsset{width=10cm,compat=1.9}

% disponi alberi
\usepackage{forest}

\forestset{
	rectstyle/.style={
		for tree={rectangle,draw,font=\large\sffamily}
	},
	roundstyle/.style={
		for tree={circle,draw,font=\large}
	}
}

% disponi algoritmi
\usepackage{algorithm}
\usepackage{algorithmic}
\makeatletter
\renewcommand{\ALG@name}{Algorithm}
\makeatother

% disponi numeri di pagina
\usepackage{fancyhdr}
\fancyhf{} 
\fancyfoot[L]{\sffamily{\thepage}}

\makeatletter
\fancyhead[L]{\raisebox{1ex}[0pt][0pt]{\sffamily{\@title \ \@date}}} 
\fancyhead[R]{\raisebox{1ex}[0pt][0pt]{\sffamily{\@author}}}
\makeatother

\begin{document}
% sezione (data)
\section{Lesson 07-10-26}

% stili pagina
\thispagestyle{empty}
\pagestyle{fancy}

% testo

\subsubsection{NAND memory}
In the last lesson we saw \textit{NOR memories}. Now we look at the \textbf{NAND} structure.

\begin{center}
	\includegraphics[scale=0.22]{../figures/nand_rom.png}
\end{center}

All bit lines are normally pulled-up to $V_{dd}$: the output is still logic 0 because of the inverter.

Now, we knot that if all transistors are at gate low, we don't connect them to ground: it follows that the bit line stays high and the output low (logic 0 at default).

The control signals for the word lines ($WL_i$) are active low.
This means that every word line is at 1 but the selected one (which is at 0).

Now:
\begin{itemize}
	\item If there is no short circuit between the drain and source of the selected transistor, we don't get a short to ground, and the line stays high: so the output stays low;
	\item If there is a short circuit, on the other hand, we short to ground, meaning the output stays high: we encode a logic 1 by shorting a transistor.
\end{itemize}

\begin{center}
	\includegraphics[scale=0.22]{../figures/nand_wl.png}
\end{center}

Clearly, the mask topology of the NAND topology is more complicated.
We have the delay of the address decoder, given by:
\begin{enumerate}
\item Delay of the address lines inverters , which are driving half of the NAND gates, and have large parasitic capacitance.
\item Delay of the NAND gates, with large parasitic resistance when all pull-down transistors are ON;
\item Lastly there's the delay of the Word Line inverters, driving all transistors in the bit cells in a row, with large parasitic capacitance and resistance as the word Line is long.
\end{enumerate}

This tradeoff comes at the gain of higher \textit{densities} overall, meaning NAND memories are usually less performant but smaller.

\subsubsection{Memory timings}

Let's look at the simulated delay, specifically of the address decoder.

\begin{center}
	\includegraphics[scale=0.22]{../figures/mem_timings.png}
\end{center}

We see the 3 delay source cited above ($t_{inv_1}$, $t_{NAND}$, and $t_{inv_2}$).
The total delay is:
$$
t_{delay} = t_{inv_1} + t_{NAND} + t_{inv_2}
$$

Note that this is the example with 64 bit cells in a row:
the delay for the first bit cell is much
lower than the delay of the last bit cell
We must always consider the worst case!

\par\smallskip

We must also look at the performance of NOR and NAND technology on bit line discharging.
The bit line has a certain parasitic capacitance $C_L$.
When the memory is idle, the bit line is pulled-up to $V_{dd}$, and $C_L$ charges up to it.
When a bit is read, we have to discharge $C_L$ down to 0.

\begin{itemize}
	\item In a \textbf{NOR} structure, the parasitic capacitance dominates.
	The resistance of the bit line is small (metal), and all capacitances are in parallel. This means that:
	$$
	\tau = R_{DSon} \cdot C_L = R_{DSon} \cdot C_{drain} \cdot n_{tr}
	$$
	where $n_{tr}$ is the number of transistors;
	\item In NAND topology the parasitic resistance dominates: it is the channel resistance of all transistors in the series.
	The structure can be studied using the \textit{Elmore} delay, and we get that:
	$$
	\tau = \frac{RC}{2}, \quad R = n_{tr} \cdot R_{DSon}, \quad C = n_{tr} \cdot C_{drain}
	$$
\end{itemize}

\par\smallskip

Latly, we show a picture of the general timing parameters for a memory access (a read, as we're talking EPROM memories) operation.

\begin{center}
	\includegraphics[scale=0.22]{../figures/mem_timings1.png}
\end{center}

From when address is stable and $\overline{CS}$ (active low) is low, we note the following delays:
\begin{itemize}
	\item $t_{OE}$ is the \textbf{minimum Output Enable} time, which is the time before the output takes some value, but is not yet stable. After $t_{OE}$ the value on DataOut might be wrong, but DataOut is not in high impedance anymore;
	\item $t_{ACS}$ is the \textbf{maxmimum Chip Select Access} time, after which the DataOut will be stable;
	\item $t_{OD}$ is the \textbf{maximum Output Disable} time, after which the output returns to high impedance.
\end{itemize}

We then see that the time we have to wait before having a result from memory is in between $t_{OE}$ and $t_{ACS}$.

\par\smallskip

Let's see what happens when changing the address.

\begin{center}
	\includegraphics[scale=0.22]{../figures/mem_timings2.png}
\end{center}

From when the address is changed, we note the following delays:
\begin{itemize}
	\item $t_{AA}$ is the \textbf{maximum Address Access} time, after which DataOut takes the value pointed to by that address;
	\item $t_{OH}$ is the \textbf{minimum Output Hold} time, after which the DataOut goes back to being undetermined after address changes;
	\item $t_{cyc}$ is the \textbf{cycle} time, for which addresses should at least be stable.
\end{itemize}

We then see that, as long as we keep addresses stable for at least $t_{cyc}$, the data will be available at most after $t_{AA}$, and at least for $t_{OH}$ after the address change.

\subsection{One Time Programmable Read Only Memories}

We now look at \textbf{OTP-ROM} (\textit{One Time Programmable Read Only Memories}).
These are based on \textbf{fuses} or \textit{anti-fuses} (which blow to become shorts), which get blown in order to encode data.
Of course this process is irreversible and means writing operations are permanent.

\begin{center}
	\includegraphics[scale=0.22]{../figures/otprom.png}
\end{center}

As in the standard NOR structure, all bit cells have a transistor, with the drain of each transistor is connected to the bit line through a fuse.

All fuses are short circuit after fabrication, meaning all transistors are connected to the bit lines: the initial content of the memory is all 1s.

Blowing a fuse turns a logic 1 into a logic 0, with he corresponding transistor being disconnected from the bit line.

As we said, the fuse cannot be restored, the operation is one time and we cannot turn a logic 0 back to a logic 1.

\par\smallskip

To correctly program an OTPROM the procedure is to:
\begin{itemize}
	\item Enable the word line of the bit we want to program;
	\item Select the bit line of the bit we want to progam;
	\item Drive the bit line at $V_p >> V_{dd}$, high enough to blow the fuse.
\end{itemize}

We call the mode in which $V_p$ is feed to the memory \textit{programming mode}.
In programming mode the data bus is an input which selects the bit lines to be connected to $V_p$.
$V_p$ itself is provided through a dedicated memory pin.

\subsubsection{Performance of OTPROMs}

First off, \textbf{writing} an OTPROM takes considerably (orders of magnitude) higher time than reading it.
This is fine, as writing an OTPROM is as we said an one-time only operation.

As for other characteristics, we see that a \textbf{read} operation is mostly identical in performance to Mask Programmable ROMs.
The disadvantages are higher resistances of the fuse compared to a short circuit, which get a bit slower compared to a Mask Programmable ROM.

OTPROMs are also \textit{unntestable}: they need to be produced in small size to have a good yield.

The number of defective ICs in a wafer does not depend on the IC size, but smaller ICs of course have less chance to be detective.
This means that usually there are more small ICs in a wafer than large ICs.
Note that the yield $y$ is roughly given by:
$$
y = \frac{n_{ICok}}{n_{ICtot}} = 1 - \frac{n_{ICdefective}}{n_{ICtot}}
$$

\subsection{Programmable Read Only Memories}

OTPROMs are suboptimal as they can be written \textit{only once}. 
Because of this, more advanced solutions which allow memory to be rewritten (that is, \textit{reprogrammed}), a number of times.
We call these \textbf{PROM}s (\textit{Programmable Read Only Memories}).

\subsubsection{Floating-gate transistors}

The core element of a PROM is a \textit{floating-gate} transistor, that is a MOS transistor which can have its gate biased by applying a high voltage to it.
One example of such a transistor is the \textbf{FAMOS} (\textit{Floating gate Avalanche MOS}), which is actually not famous at all.

They are similar to a MOS transistor, with an dditional polysilicon layer between two thin $SiO2$ layers below the gate.
This means the gate is \textit{floating}.
Charges trapped in the floating gate change the threshold voltage $V_T$, where 
\begin{itemize}
	\item 
Negative charges increase the threshold;
	\item
Positive charges decrease the threshold.
\end{itemize}

\newpage
\textbf{\textsf{Programming}}

The floating gate is completely isolated, and to program it we use \textit{hot electron programming}:
\begin{itemize}
	\item By applying a high programming voltage $V_p$ to the gate, a rich channel forms (as the gate is pulled high). This is not enough to lower $V_T$!
	\item By applying a high programming voltage $V_p'$ to the drain-source channel, electrons in the channel are subject to strong electrical fields.
Due to strong vertical electrical field from the gate, the high energy electrons may overcome the $SiO2$ barrier, reach the floating gate. The result is that the gate is pulled low and $V_T$ increases.
\end{itemize}

\par\smallskip
\textbf{\textsf{Performance}}

When programming a FAMOS transistor, as the threshold voltage $V_T$ increases, the channel is less rich and the vertical electrical field is less strong.
This clearly means that the number of electrons with enough energy to jump in the floating gate decreases (less electrical field), and process is \textit{self regulating} (!)
The nice property of this system is that the final threshold voltage is quite consistent among all transistors, regardless of the length of the process.

We then have that $V_T$ grows above $V_{dd}$ by a certain margin.
By playing with doping levels, densities, geometries, we can change that margin.

\begin{center}
	\includegraphics[scale=0.22]{../figures/famos.png}
\end{center}

The \textit{disadvantages} of FAMOS transistors are:
\begin{itemize}
	\item The process takes time and requires high voltages
	\item How to generate those voltages in the first place: provided externally using pins? Generated internally from the standard voltage supply?
	\item Damages to the lattice: we cannot apply programming too many times. 1000 – 10000 program ing cycles seems to be the sweet spot in most commercial applications;
	\item Leakage current removes electrons from the floating gate: the information is lost over time, on a timescale of sveral tens of years.
\end{itemize}

\par\smallskip
\textbf{\textsf{Erasing}}

\textit{Erasing} FAMOS transistors is the most interesting part.
Electrons cannot be removed electrically from the floating gate: energy has to provided to the electrons via an \textit{UV lamp}.

As high energy photons hit electrons, the electrons jump back to the substrate and diffuse.
After a sufficiently long amount of time (minutes) all electrons are removed, and the floating gate is left unbiased.

For this reason we need a window on the integrated circuit (which we'll tape off when the circuit is at work to avoid accidental exposure to UV light).
We also cannot target a \textit{single} transistor: all transistors are erased at the same time.
This makes PROM erasure an usually out of the circuit operation, as with programming, done with external hardware in a lab. 

\subsubsection{Programmable ROM}

Let's look at a NOR structure \textbf{PROM}.

\begin{center}
	\includegraphics[scale=0.22]{../figures/prom.png}
\end{center}

The funky transistor symbol with the triple gate line represents a FAMOS transistor (with the middle line being the floating gate).

All bit cells have a FAMOS transistor.
The drain of each transistor is connected to the bit line.
All FAMOS transistors are not programmed after fabrication, so they behave normally.
The initial content of the memory, as a result, is all 1s, and programming a transistor (raising the $V_T$ higher, disconnecting it) turns a logic 1 into a logic 0.

As we said, the memory can be reset to the initial state by shining UV rays through a window.
All programmed transistors are erased, and all non programmed transistors are not affected: the memory contents return to all 1s. 

To correctly program a PROM we proceed similarly to OTPROMs:
\begin{itemize}
	\item Enable the word line of the bit we want to program;
	\item Select the bit line of the bit we want to program;
	\item The selected word line is driven to the voltage $V_p > V_{dd}$;
	\item The selected bit lines (possibly more than one for the same word) are driven to the voltage $V_p' > V_{dd}$.
\end{itemize}

To avoid impacting other components, other rows are at ground, and bits not to be programmed on this row have their drains at $V_{dd}$.

\subsubsection{Performance of PROMs}

For the read operation, PROMs are identical to Mask Programmable ROM and OTPROMs.

For the \textit{write} operation, that is \textit{programming},
moving negative charges to the floating gate
makes the transistor a permanent open circuit:
\begin{itemize}
	\item Turns a 1 into a 0 (if the bit line is inverted);
	\item Turns a 0 into a 1 (if the bit line is affirmed).
\end{itemize}

It works at the word level too.

For erasing, we emove the negative charges from the floating gate.
This restores the normal transistor behavior, but applies to the entire memory at once.

This means PROMs have almost completely superseded OTPROMs:
\begin{itemize}
\item Much more flexible;
\item Cheaper during prototyping;
\item Can be reused;
\item Still, not suitable for data that is frequently changed.
\end{itemize}

\subsubsection{FLOTOX transistors}

Another variant of floating-gate transistor is the \textbf{FLOTOX} (\textit{FLOating gate Tunneling OXide transistor}). Similarly to FAMOS, it has a floating gate.

The FLOTOX differs from the FAMOS in its erasure properties: it is \textit{electrically erasable}.

This is achieved through the \textit{tunneling effect}.
The tunneling effect is  a quantum mechanical phenomenon, where a particle (electron) can pass a potential barrier even if it has insufficient energy.

A thinner oxide layer on the drain (source) side makes tunneling easier.
Electrons come from the drain (source) region rather than from the substrate or channel.
An higher electrical field is necessary for this.

\par\smallskip
\textbf{\textsf{Programming}}

The floating gate is completely isolated, and to program it we use
\textit{Fowler--Nordheim tunneling}:
\begin{itemize}
    \item By applying a high programming voltage $V_p$ to the control
    gate, a strong electric field is established across the thin
    tunnel oxide.

    \item The strong electric field allows electrons to tunnel through
    the thin $SiO_2$ barrier from the substrate onto the floating gate.

    \item The floating gate accumulates negative charge. As a result,
    the threshold voltage increases:
    \[
        Q_{FG}<0 \quad\Rightarrow\quad V_T \uparrow.
    \]
\end{itemize}

To erase the cell, the polarity of the applied voltage is reversed:
\begin{itemize}
    \item A high voltage of opposite polarity is applied to the control
    gate, producing a strong electric field across the tunnel oxide.

    \item Electrons tunnel from the floating gate through the thin
    $SiO_2$ layer into the substrate.

    \item The floating gate loses its negative charge, causing the
    threshold voltage to decrease:
    \[
        Q_{FG}\rightarrow 0 \quad\Rightarrow\quad V_T \downarrow.
    \]
\end{itemize}

\par\smallskip
\textbf{\textsf{Performance}}

Let's make a comparison of programming mechanisms:
\begin{itemize}

	\item Hot Electron Programming is faster than Tunneling, which may increase performance, but performance may be limited by other issues.

	\item Tunneling causes less damages to the lattice than Hot Electron Programming, leading to higher life-time of the device, and many more programming cycles;

	\item Hot Electron Programming is Self Regulated, Tunneling is not;

	\item Tunneling requires careful timing to achieve good consistency:
	the resulting distribution of $V_T$ after programming has a larger standard deviation,
	and requires larger margins during operations;

	\item FLOTOX transistors are larger and harder to fabricate: lower yield, higher costs.
\end{itemize}

\par\smallskip
\textbf{\textsf{Addenda on programming}}

There is a famous issue with FLOTOX programming.
The straightforward implementation (but it doesn't work!) is:
\begin{itemize}
	\item The drain of each transistor is connected to the bit line;
	\item The source of each transistor is grounded.
\end{itemize}
This has a problem: programming a transistor erases another one.

\begin{itemize}
    \item To program $T_1$:
    $
        \text{Gate: } WL_0 = 12\,V,
        \qquad
        \text{Drain: } BL_1 = 0\,V
    $

    \item To avoid programming $T_2$:
    $
        WL_1 = 0\,V
    $

    \item To avoid programming $T_3$:
    $
        BL_0 = 12\,V,
        \qquad
        V_{GD} = 0
    $

    \item However, $T_4$ is erased:
    $
        WL_1 = 0\,V,
        \qquad
        BL_0 = 12\,V
    $

\end{itemize}

The solution is to use an access transistor for each bit cell.

\begin{center}
	\includegraphics[scale=0.22]{../figures/flotox_prog.png}
\end{center}

\subsubsection{Electrically Erasable Programmable ROM}

We use FLOTOX transistors to realize \textbf{EEPROM}s (\textit{Electrically Erasable Programmable ROMs}). 

\begin{center}
	\includegraphics[scale=0.22]{../figures/eeprom.png}
\end{center}

All bit cells have a FLOTOX transistor, with thin oxide on the drain side.

All bit cells also have an access transistor;
drain connected to the bit bine, source connected to the FLOTOX,
gate driven by the address decoder (via dedicated $AWL_i$ lines). 

Access transistors are activated when the word is selected.

Out of the factory, all transistors behave normally, and the initial content of the memory is all 1s.
\begin{itemize}
	\item Programming a transistor turns a logic 1 into a logic 0;
	\item Erasing a transistor turns a logic 0 back into a logic 1.
\end{itemize}

This happens \textit{electrically} and \textit{in circuit}.

Let's look at how the main operations are carried out:
\begin{itemize}
	\item For \textbf{programming}, we select the word by raising the word line, and the data in the word by raising the D line (which as always acts as an input for programming).
		
		We then enter the \textit{programming} mode by enabling $AWL$ on the word line, that is enabling the access transistors. We keep these OFF on other lines to avoid affecting other FLOTOX transistors.

		At this point, to program we raise the word line (that is the gate) to $V_p$, and the selected bit lines (that is the drains) to $0 \, \matrhm{V}$. The rest of the lines are kept at $V_p$; 

	\item For \textbf{erasing}, we select the word by raising the word line, and the data (to erase) in the word by raising the D line.

		We then enter the \textit{erasing} mode by enabling $AWL$ on the word line, that is enabling the access transistors. We keep these OFF on other lines to avoid affecting other FLOTOX transistors.

		At this point, to erase we lower the word line (that is the gate) to $0 \, \mathrm{V}$, and the selected bit lines (that is the drains) to $V_p$. The rest of the lines are kept at $0 \, \mathrm{V}$.
\end{itemize}

\subsubsection{Performance of EEPROMs}

The read operation is identical to Mask Programmable ROM, OTPROM and EPROMs (which we called PROMs so far).
The eternal interface is to:
\begin{itemize}
	\item Provide programming voltages;
	\item Signal to select programming;
	\item Signal to select erasing.
\end{itemize}
External interface often simplified by providing internal programming and erasing mechanisms.

They are widely used, and much more flexible than EPROMs, as they can be erased by a micro-processor in circuit.
On the other hand, they are expensive, and the bit cell is larger than in an EPROM.

\par\smallskip

\textit{Writeability} is very good as it is:
\begin{itemize}
	\item In circut;
	\item Repeatable, thousands of times.
\end{itemize}

We also have \textit{erasure} in circuit, with single bit access both in writing and erasing.

The storage permanence is very high, and the NREs are kept low (though they are higher than EPROMs).

\end{document}
